\Askhsh{27320}{4}
\begin{ylh}{1}
    \item 2.6 Συνέπειες του Θεωρήματος Μέσης Τιμής
    \item 2.7 Τοπικά ακρότατα συνάρτησης
    \item 2.8 Κυρτότητα - Σημεία καμπής συνάρτησης.
\end{ylh}
\begin{wrapfigure}[15]{r}{8cm} % Adjust the number of lines to wrap [15] as needed
\centering
\begin{tikzpicture}
\begin{axis}[
    axis lines = middle,
    xlabel = {$x$},
    ylabel = {$y$},
    xmin = 0, xmax = 5,
    ymin = -2.5, ymax = 2.5,
    xtick = {1},
    ytick = {}, % No y-ticks shown explicitly in the image except for the origin.
    grid = none, % No grid shown in the image
    samples = 100,
    width = 7cm, height = 7cm,
    tick label style = {font=\scriptsize},
    every axis x label/.style = {at={(current axis.right of origin)}, anchor=west},
    every axis y label/.style = {at={(current axis.above origin)}, anchor=south},
    clip = false, % Allows labels outside the axis box
]
    % Drawing the curve for f'(x). A logarithmic function fits the description.
    \addplot[very thick, maincolor, domain=0.01:4] {ln(x)};
    \node[maincolor, above right] at (axis cs:4, {ln(4)}) {$C_{f'}$}; % Label for the curve
    
    % Point A(1,0)
    \coordinate (A_point) at (axis cs:1,0); % Define point A
    \fill (A_point) circle (2pt); % Draw point A
    \node[below right] at (A_point) {$A(1,0)$}; % Label for point A
    
    % Origin O
    \node[below left] at (axis cs:0,0) {O}; % Label for origin
\end{axis}
\end{tikzpicture}
\end{wrapfigure}
Στο παρακάτω σχήμα δίνεται στο $(0, +\infty)$ η γραφική παράσταση της παραγώγου $f'$ μιας συνάρτησης $f$ με πεδίο ορισμού το $(0, +\infty)$. Δίνεται επίσης ότι η $f'$ είναι συνεχής και γνησίως αύξουσα συνάρτηση στο $(0, +\infty)$ με $\lim_{x \to +\infty} f'(x) = +\infty$.

\begin{alist}
\item Να βρείτε τα διαστήματα μονοτονίας και τα τοπικά ακρότατα της συνάρτησης $f$. \monades{09}
\item Ένας μαθητής ισχυρίζεται ότι:
    \begin{rlist}
    \item «Η γραφική παράσταση της $f$ δέχεται οριζόντια εφαπτομένη στο σημείο με τετμημένη $1$».
    \item «Υπάρχει μοναδικό $\kappa \in (0, +\infty)$ τέτοιο, ώστε ο συντελεστής διεύθυνσης της εφαπτομένης της $C_f$ στο σημείο $M(\kappa, f(\kappa))$ να ισούται με $2$».
    \end{rlist}
    Ποιοί από τους παραπάνω ισχυρισμούς του μαθητή είναι σωστοί; Να δικαιολογήσετε τις απαντήσεις σας. \monades{10}
\item Τι μπορούμε να πούμε για την κυρτότητα της $f$ στο πεδίο ορισμού της; Να δικαιολογήσετε την όποια απάντησή σας. \monades{06}
\end{alist}